\documentclass[11pt]{amsart}
\usepackage[margin=1.1in]{geometry}
\usepackage{amsmath,amssymb,amsthm,mathtools}
\usepackage{booktabs,array}
\usepackage[colorlinks=true,linkcolor=black,citecolor=black,urlcolor=blue]{hyperref}

\newtheorem{theorem}{Theorem}[section]
\newtheorem{proposition}[theorem]{Proposition}
\newtheorem{lemma}[theorem]{Lemma}
\newtheorem{corollary}[theorem]{Corollary}
\newtheorem*{mainthm}{Main Theorem}
\theoremstyle{definition}
\newtheorem{definition}[theorem]{Definition}
\theoremstyle{remark}
\newtheorem{remark}[theorem]{Remark}

\numberwithin{equation}{section}

\newcommand{\e}{\mathbf{e}}
\newcommand{\Z}{\mathbb{Z}}
\newcommand{\Q}{\mathbb{Q}}
\newcommand{\C}{\mathbb{C}}
\newcommand{\re}{\operatorname{Re}}
\newcommand{\im}{\operatorname{Im}}
\newcommand{\THq}{\mathcal{T}}
\newcommand{\SL}{\mathrm{SL}_2(\Z)}
\newcommand{\Ahat}{\widehat{A}}
\newcommand{\Bhat}{\widehat{B}}
\newcommand{\fr}[1]{\{#1\}}

\begin{document}

\title[Sign laws for mock theta functions]{Sign laws for the coefficients of second- and eighth-order mock theta functions}

\author{Ramanuja}
\address{Kachhubi, P.O.\ Tamuria, Lakhnaur Block, Madhubani, Bihar 847410, India}
\email{\href{mailto:ramanujamath@yahoo.com}{ramanujamath@yahoo.com}}

\subjclass[2020]{11P82, 11F37, 11F27, 11P55}
\keywords{Mock theta functions, Appell--Lerch sums, circle method, sign patterns, Siegel functions}

\date{November 25, 2026}

\begin{abstract}
We prove periodic sign laws for the coefficients of five mock theta functions: the eighth-order functions $S_0,S_1,U_0,U_1$ of Gordon and McIntosh and the second-order function $\mu$.
For $S_0,S_1,\mu,U_1$ the sign of the $n$-th coefficient depends only on $n \bmod 4$, and for $U_0$ only on $n\bmod 8$, apart from an explicit finite list of exceptions.
We obtain asymptotic formulas with exact constants, for instance
$c_{S_0}(n)\sim \cos\!\big(\tfrac{9\pi}{16}-\tfrac{\pi n}{2}\big)\cosh\!\big(\tfrac{\pi}{4}\sqrt{n-\tfrac1{16}}\big)\big/\sqrt{n-\tfrac1{16}}$,
together with a fully explicit error bound from the circle method.
The bound forces the sign law for $n\ge N_0$ with $N_0\le 11{,}231$, and an exact computation covers the remaining $n$.
The proof writes each function as Appell--Lerch sums plus a theta quotient, shows that the dominant growth always comes from a pure theta quotient, and bounds every theta factor uniformly on every Farey arc by means of Siegel functions.
Along the way we prove the identity $U_0+2U_1=J_2^8/(J_1^3J_4^4)$.
\end{abstract}

\maketitle

%=====================================================================
\section{Introduction}

The coefficients of several of Ramanujan's mock theta functions show sign patterns of striking regularity.
The prototype is Ramanujan's claim about the third-order function $f(q)$, whose coefficients alternate in sign; their asymptotics were studied by Dragonette~\cite{Dr} and Andrews~\cite{An}, and an exact formula was proved by Bringmann and Ono~\cite{BO}.
Recently Dastidar~\cite{Da} proved that the coefficients $r(n)$ of the third-order function $\rho(q)$ satisfy
$r(3m)>0$, $r(3m+1)\le 0$, $r(3m+2)\le 0$, with equality only for $n\in\{2,4,8,11,20\}$, and he proposed a systematic study of such sign laws for all mock theta functions.

In this paper we carry out that study for the eighth-order functions of Gordon and McIntosh~\cite{GM} and the second-order function $\mu$ of McIntosh~\cite{Mc}.
For $|q|<1$ we use the $q$-Pochhammer symbol $(x;q)_n=\prod_{j=0}^{n-1}(1-xq^j)$ and define
\begin{align}
S_0(q)&=\sum_{n\ge0}\frac{q^{n^2}(-q;q^2)_n}{(-q^2;q^2)_n}, &
S_1(q)&=\sum_{n\ge0}\frac{q^{n(n+2)}(-q;q^2)_n}{(-q^2;q^2)_n}, \label{eq:S}\\
U_0(q)&=\sum_{n\ge0}\frac{q^{n^2}(-q;q^2)_n}{(-q^4;q^4)_n}, &
U_1(q)&=\sum_{n\ge0}\frac{q^{(n+1)^2}(-q;q^2)_n}{(-q^2;q^4)_{n+1}}, \label{eq:U}\\
\mu(q)&=\sum_{n\ge0}\frac{(-1)^nq^{n^2}(q;q^2)_n}{(-q^2;q^2)_n^2}. \label{eq:mu}
\end{align}
For $f\in\{S_0,S_1,\mu,U_0,U_1\}$ write $f(q)=\sum_{n\ge0}c_f(n)q^n$; all $c_f(n)$ are integers.

\begin{mainthm}
For each $f$ in Table~\ref{tab:main} and every $n\ge0$, the sign of $c_f(n)$ is the one listed for the residue class of $n$, except exactly at the listed values of $n$, where $c_f(n)$ takes the listed value.
\end{mainthm}

\begin{table}[ht]
\centering
\renewcommand{\arraystretch}{1.25}
\small
\begin{tabular}{@{}l>{\raggedright\arraybackslash}p{4.9cm}>{\raggedright\arraybackslash}p{5.4cm}@{}}
\toprule
$f$ & sign of $c_f(n)$ & exceptions \\
\midrule
$S_0$ & $-,+,+,-$ for $n\equiv 0,1,2,3 \pmod 4$ & $c(0)=1$; \ $c(n)=0$ for $n=4,6,8,12,14,20,30$\\
$S_1$ & $+,-,-,+$ for $n\equiv 0,1,2,3 \pmod 4$ & $c(n)=0$ for $n=1,2,9$\\
$\mu$ & $-,-,+,+$ for $n\equiv 0,1,2,3 \pmod 4$ & $c(0)=1$\\
$U_1$ & $+,+,-,-$ for $n\equiv 0,1,2,3 \pmod 4$ & $c(n)=0$ for $n=0,2$\\
$U_0$ & $-,+,+,-,+,-,-,+$ for $n\equiv 0,\dots,7 \pmod 8$ & $c(0)=1$; \ $c(n)=0$ for $n=3,5,8,12,16,19,24,28,40,60$\\
\bottomrule
\end{tabular}
\caption{The sign laws.}\label{tab:main}
\end{table}

The Main Theorem follows from asymptotic formulas with exact constants (Theorem~\ref{thm:asym}) combined with an explicit error bound (Theorem~\ref{thm:explicit}).
For example,
\begin{equation}\label{eq:intro-asym}
c_{S_0}(n)=\cos\Big(\frac{9\pi}{16}-\frac{\pi n}{2}\Big)\frac{\cosh\big(\frac{\pi}{4}\sqrt{n-\frac1{16}}\big)}{\sqrt{n-\frac1{16}}}+O\big(n^{3/4}e^{\pi\sqrt n/8}\big).
\end{equation}
In every case the trigonometric factor depends only on $n$ modulo $4$ or $8$ and never vanishes (Lemma~\ref{lem:trig}), so the sign law holds as soon as the error is smaller than the main term.
We make the error bound completely explicit; it suffices for $n\ge N_0$ with $N_0=11231,\,10652,\,976,\,493,\,5180$ for $S_0,S_1,\mu,U_1,U_0$ respectively.
The finitely many remaining $n$ are checked by exact integer computation.

\subsection*{Method}
Hickerson and Mortenson~\cite{HM} expressed every classical mock theta function through Appell--Lerch sums $m(x,q,z)$ and theta quotients.
For our five functions we show that the dominant exponential growth, located at $q=\pm i$ for $S_0,S_1,\mu,U_1$ and at the primitive eighth roots of unity for $U_0$, comes entirely from a pure theta quotient, while the Appell--Lerch part is bounded there.
The main terms are therefore given by modular transformation of theta functions, with constants that we compute exactly.
For the error we need upper bounds for the Appell--Lerch parts and two-sided bounds for theta functions on all Farey arcs.
We obtain both uniformly by writing each theta factor as a Siegel function times the Dedekind eta function and using their transformation laws.
A new identity, $U_0+2U_1=J_2^8/(J_1^3J_4^4)$ (Proposition~\ref{prop:TH}), reduces the three functions of base $q^4$ to a single eta quotient.

\subsection*{Computations}
The proof uses two computer calculations, both described in Appendix~\ref{app:code}: the exact integer coefficients for $n<N_0$, and the evaluation of finitely many explicit constants (convergent infinite products) in multiprecision arithmetic with outward rounding.
The source code is provided as ancillary files and at \url{https://github.com/ramanujamath/mock-theta-sign-laws}.

\subsection*{Notation}
Throughout, $\e(x)=e^{2\pi i x}$, $J_m=(q^m;q^m)_\infty$, and $\fr{x}=x-\lfloor x\rfloor$ denotes the fractional part.
For a cusp $h/k$ with $\gcd(h,k)=1$ we write $q=\e(h/k)\,e^{-z}$ with $\re z>0$.
The second Bernoulli polynomial is $B_2(x)=x^2-x+\frac16$.

%=====================================================================
\section{Appell--Lerch representations}\label{sec:AL}

The Jacobi theta function is
\[
\Theta(x;q)=\sum_{n\in\Z}(-1)^nq^{n(n-1)/2}x^n=(x;q)_\infty(q/x;q)_\infty(q;q)_\infty \qquad (x\neq0,\ |q|<1),
\]
the second expression being the Jacobi triple product.
It satisfies $\Theta(x^{-1};q)=-x^{-1}\Theta(x;q)$.
Following~\cite{HM}, the Appell--Lerch sum is defined for $x,z\neq0$ with $z\notin q^{\Z}$ and $xz\notin q^{\Z}$ by
\begin{equation}\label{eq:m}
m(x,q,z)=\frac{1}{\Theta(z;q)}\sum_{r\in\Z}\frac{(-1)^rq^{r(r-1)/2}z^r}{1-q^{r-1}xz}.
\end{equation}
We use three properties from~\cite[Prop.~2.1, Cor.~2.4, Thm.~2.3]{HM}:
\begin{align}
m(x,q,z)&=m(x,q,qz), \label{eq:m-shift}\\
m(x,q,z)&=m(x,q,x^{-1}z^{-1}), \label{eq:m-sym}\\
m(x,q,z_1)-m(x,q,z_0)&=\frac{z_0J_1^3\,\Theta(z_1/z_0;q)\,\Theta(xz_0z_1;q)}{\Theta(z_0;q)\Theta(z_1;q)\Theta(xz_0;q)\Theta(xz_1;q)}, \label{eq:m-change}
\end{align}
where in \eqref{eq:m-change} $J_1=(q;q)_\infty$ is taken with respect to the base $q$ of $m$.

\begin{proposition}[Hickerson--Mortenson {\cite[\S4]{HM}}]\label{prop:HM}
As formal power series and as analytic functions on $|q|<1$:
\begin{enumerate}
\item $\displaystyle S_0(q)=2m(-q^3,q^8,-1)+\frac{q\,\Theta(-q;q^8)\Theta(q^2;q^8)^2}{\Theta(q^3;q^8)^2}$\newline and $S_0(q)=m(-q^3,q^8,-q^2)+m(-q^3,q^8,-q^6)$;
\item $\displaystyle S_1(q)=-2q^{-1}m(-q,q^8,-1)+\frac{\Theta(-q^3;q^8)\Theta(q^2;q^8)^2}{q\,\Theta(q;q^8)^2}$\newline and $S_1(q)=-q^{-1}m(-q,q^8,-q^2)-q^{-1}m(-q,q^8,-q^6)$;
\item $U_0(q)=2m(-q,q^4,-1)$ and $U_1(q)=-m(-q,q^4,-q^2)$;
\item $\mu(q)=2m(-q,q^4,-1)+2m(-q,q^4,q)$.
\end{enumerate}
\end{proposition}

For $S_0$ and $S_1$ we call the first expression \emph{form A} and the second \emph{form B}.

\begin{corollary}\label{cor:U}
$U_0(q)=2m(-q,q^4,q^3)$ and $U_1(q)=-m(-q,q^4,q)$.
\end{corollary}

\begin{proof}
Apply \eqref{eq:m-sym} with $x=-q$ and base $q^4$: the parameter $z=-1$ becomes $x^{-1}z^{-1}=q^{-1}$, and $z=-q^2$ becomes $q^{-3}$.
Then \eqref{eq:m-shift} with base $q^4$ replaces $q^{-1}$ by $q^3$ and $q^{-3}$ by $q$.
\end{proof}

\begin{proposition}\label{prop:TH}
Let $\THq(q)=J_2^8/(J_1^3J_4^4)$.
Then $U_0+2U_1=\THq$.
Consequently
\[
\mu=2U_0-\THq,\qquad U_1=\tfrac12(\THq-U_0),\qquad U_0=\THq-2U_1.
\]
\end{proposition}

\begin{proof}
By Proposition~\ref{prop:HM}(3), $U_0+2U_1=2\big[m(-q,q^4,-1)-m(-q,q^4,-q^2)\big]$.
Apply \eqref{eq:m-change} with base $Q=q^4$, $x=-q$, $z_1=-1$, $z_0=-q^2$.
Then $z_1/z_0=q^{-2}$, $xz_0z_1=-q^3$, $xz_0=q^3$ and $xz_1=q$, so
\[
U_0+2U_1=-2q^2J_4^3\cdot\frac{\Theta(q^{-2};q^4)}{\Theta(-q^2;q^4)\,\Theta(-1;q^4)}\cdot\frac{\Theta(-q^3;q^4)}{\Theta(q^3;q^4)\,\Theta(q;q^4)}.
\]
By the triple product and $\Theta(q^{-2};q^4)=-q^{-2}\Theta(q^2;q^4)$,
\begin{gather*}
\Theta(q^2;q^4)=(q^2;q^4)_\infty^2J_4,\qquad \Theta(-q^3;q^4)=(-q;q^2)_\infty J_4,\qquad \Theta(-q^2;q^4)=(-q^2;q^4)_\infty^2J_4,\\
\Theta(-1;q^4)=2(-q^4;q^4)_\infty^2J_4,\qquad \Theta(q;q^4)=\Theta(q^3;q^4)=(q;q^2)_\infty J_4.
\end{gather*}
Substituting, and using $(-q^2;q^4)_\infty(-q^4;q^4)_\infty=(-q^2;q^2)_\infty$,
\[
U_0+2U_1=\frac{J_4\,(q^2;q^4)_\infty^2\,(-q;q^2)_\infty}{(-q^2;q^2)_\infty^2\,(q;q^2)_\infty^2}.
\]
Finally $(q;q^2)_\infty=J_1/J_2$, $(q^2;q^4)_\infty=J_2/J_4$, $(-q;q^2)_\infty=J_2^2/(J_1J_4)$ and $(-q^2;q^2)_\infty=J_4/J_2$, which give $J_2^8/(J_1^3J_4^4)$.
For the consequences: by Proposition~\ref{prop:HM}(4) and Corollary~\ref{cor:U}, $\mu=U_0-2U_1$, so $\mu=U_0-(\THq-U_0)=2U_0-\THq$.
\end{proof}

%=====================================================================
\section{Theta functions on Farey arcs}\label{sec:theta}

Every theta factor appearing in Section~\ref{sec:AL} has the form $\Theta(s\,q^a;q^M)$ with $M\in\{4,8\}$, $0\le a<M$ and $s=\e(\sigma)$, $\sigma\in\{0,\frac12\}$, or is an eta quotient $J_m$, $m\in\{1,2,4\}$; note $J_m=\Theta(q^m;q^{3m})$.
Write $q=\e(\tau)$ with $\im\tau>0$.

\begin{definition}
For $(A,B)\in\Q^2$ the \emph{Siegel function} is
\[
g_{(A,B)}(T)=-\mathcal Q^{B_2(A)/2}\,\e\big(B(A-1)/2\big)\,(1-q_z)\prod_{n\ge1}(1-\mathcal Q^nq_z)(1-\mathcal Q^n/q_z),
\]
where $\mathcal Q=\e(T)$ and $q_z=\e(AT+B)$,
and $\eta(T)=\mathcal Q^{1/24}\prod_{n\ge1}(1-\mathcal Q^n)$ is the Dedekind eta function.
\end{definition}

\begin{lemma}\label{lem:siegel}
Let $A=a/M$. Then
\[
|\Theta(s\,q^a;q^M)|=|q|^{M(-\frac18+\frac A2-\frac{A^2}2)}\;|g_{(A,\sigma)}(M\tau)|\;|\eta(M\tau)|.
\]
\end{lemma}

\begin{proof}
Compare the triple product for $\Theta(sq^a;q^M)$ with the product formulas for $g_{(A,\sigma)}(M\tau)$ and $\eta(M\tau)$.
The powers of $\mathcal Q=q^M$ differ by $\mathcal Q^{-B_2(A)/2-1/24}$, whose modulus is the stated power of $|q|$.
\end{proof}

\begin{lemma}[Transformation]\label{lem:transf}
Let $q=\e(h/k)e^{-z}$, $\re z>0$, and $M\ge1$.
Put $g=\gcd(M,k)$, $k'=k/g$, $h'=Mh/g$, choose $\gamma=\begin{psmallmatrix}\alpha&\beta\\ k'&-h'\end{psmallmatrix}\in\SL$, and let $T'=\gamma(M\tau)$.
Then
\begin{enumerate}
\item $|\eta(M\tau)|=\big(2\pi g/(Mk|z|)\big)^{1/2}\,|\eta(T')|$;
\item $|g_{(A,\sigma)}(M\tau)|=|g_{(\Ahat,\Bhat)}(T')|$ with $\Ahat=\fr{-Ah'-\sigma k'}$ and $\Bhat\in\frac1M\Z$;
\item $\rho:=|\e(T')|=\exp\!\big(-4\pi^2g^2\re(1/z)/(Mk^2)\big)$.
\end{enumerate}
\end{lemma}

\begin{proof}
We have $k'M\tau-h'=ik'Mz/(2\pi)$, so (1) is the transformation law of $\eta$ under $\gamma$ (see~\cite[Thm.~3.4]{Ap}), and (3) follows from $\im T'=\im(M\tau)/|k'M\tau-h'|^2$ and $\im(M\tau)=M\re z/(2\pi)$.
Statement (2) is the transformation law of Siegel functions, $g_{a}\circ\gamma=\varepsilon\, g_{a\gamma}$ with $|\varepsilon|=1$, together with the fact that $|g_a|$ depends only on $a\bmod\Z^2$~\cite[Ch.~2, K1--K2]{KL}.
Here $(A,\sigma)\gamma^{-1}=(-Ah'-\sigma k',\ \sigma\alpha-A\beta)$, written with respect to the action used in~\cite{KL}.
\end{proof}

\begin{lemma}[Two-sided bounds]\label{lem:bounds}
Let $(\Ahat,\Bhat)\in(\frac18\Z)^2\setminus\Z^2$ with $0\le\Ahat<1$, and let $\rho=|\e(T')|<1$.
Then
\[
\rho^{B_2(\Ahat)/2}L(\Ahat,\rho)\le|g_{(\Ahat,\Bhat)}(T')|\le\rho^{B_2(\Ahat)/2}U(\Ahat,\rho),
\]
where
\begin{align*}
U(\Ahat,\rho)&=u_0\prod_{n\ge1}(1+\rho^{n+\Ahat})(1+\rho^{n-\Ahat}),&u_0&=1+\rho^{\Ahat}\ (\Ahat>0),\ u_0=2\ (\Ahat=0),\\
L(\Ahat,\rho)&=\ell_0\prod_{n\ge1}(1-\rho^{n+\Ahat})(1-\rho^{n-\Ahat}),&\ell_0&=1-\rho^{\Ahat}\ (\Ahat>0),\ \ell_0=2\sin(\pi/8)\ (\Ahat=0).
\end{align*}
Likewise $\rho^{1/24}\prod_{n\ge1}(1-\rho^n)\le|\eta(T')|\le\rho^{1/24}\prod_{n\ge1}(1+\rho^n)$.
The functions $U$ and $\prod(1+\rho^n)$ increase with $\rho$, and $L$ and $\prod(1-\rho^n)$ decrease.
\end{lemma}

\begin{proof}
Since $|q_z|=\rho^{\Ahat}$, the bounds follow from $1-|w|\le|1-w|\le1+|w|$ applied to every factor.
If $\Ahat=0$ then $\Bhat\in\frac18\Z\setminus\Z$, and $|1-\e(\Bhat)|\ge2\sin(\pi/8)$.
All exponents $n\pm\Ahat$ are positive because $\Ahat\in\frac18\Z\cap[0,1)$.
\end{proof}

\begin{definition}[Growth exponent]\label{def:E}
Consider a product $P=\prod_j\Theta_j^{p_j}$ ($p_j\in\Z$) of factors $\Theta_j=\Theta(s_jq^{a_j};q^{M_j})$ or $\Theta_j=J_{m_j}$.
On the Farey arc of $h/k$ define
\[
E_P(h/k)=-\frac{256}{k^2}\sum_j p_j\,\frac{g_j^2}{M_j}\Big(\frac{B_2(\Ahat_j)}2+\frac1{24}\Big),
\]
where $g_j,\Ahat_j$ are as in Lemma~\ref{lem:transf} for the $j$-th factor; for $\Theta_j=J_m$ take $M_j=m$ and replace $B_2(\Ahat_j)/2+1/24$ by $1/24$.
\end{definition}

By Lemmas~\ref{lem:siegel}--\ref{lem:bounds}, $|P|$ equals
\[
\exp\Big(\frac{\pi^2}{64}E_P(h/k)\,\re(1/z)\Big)
\]
times explicit factors that are bounded above and below by products of the form $U$, $L$, $\prod(1\pm\rho^n)$ and powers of $k|z|$ and $|q|$.
The number $k^2E_P(h/k)$ depends only on $(k\bmod 2M,\ h\bmod M)$ with $M=\operatorname{lcm}(M_j)$, because $g_j$ depends on $k\bmod M_j$, while $\Ahat_j$ depends on $h'\bmod M_j$ and $k'\bmod2$.
Hence $E_P$ takes finitely many values times $k^{-2}$, and its maximum over all cusps is attained at small $k$.

\begin{proposition}[Exponents]\label{prop:exponents}
Use the following representations on the arc of $h/k$:
\begin{itemize}
\item for $S_0$ and $S_1$: form A if $k'=k/\gcd(8,k)$ is odd, form B if $k'$ is even;
\item for $\mu$: $\mu=-\THq+2U_0$, with $U_0$ written by Proposition~\ref{prop:HM}(3) or Corollary~\ref{cor:U};
\item for $U_1$: $U_1=\frac12\THq-\frac12U_0$ (same choices for $U_0$), or one of the two single-$m$ forms of $U_1$;
\item for $U_0$: one of the two single-$m$ forms of $U_0$, or $U_0=\THq-2U_1$ (single-$m$ forms for $U_1$).
\end{itemize}
In each case take the choice with the smallest maximal exponent.
For an Appell--Lerch term $c\,m(x,q^M,z)$ the relevant product is $\Theta(z;q^M)^{-1}$.
Then the largest exponent of any term is as in Table~\ref{tab:exp}.
\end{proposition}

\begin{table}[ht]
\centering
\renewcommand{\arraystretch}{1.2}
\begin{tabular}{@{}lccc@{}}
\toprule
$f$ & dominant cusps $\zeta$ & exponent at $\zeta$ & largest exponent at all other cusps, $\varepsilon_f$\\
\midrule
$S_0,S_1$ & $\pm i$ & $1$ & $1/4$\\
$\mu$ & $\pm i$ & $2$ & $1/2$\\
$U_1$ & $\pm i$ & $2$ & $2/9$\\
$U_0$ & $\e(h/8)$, $h$ odd & $1/2$ & $1/8$\\
\bottomrule
\end{tabular}
\caption{Growth exponents, in units of $\pi^2\re(1/z)/64$.}\label{tab:exp}
\end{table}

\begin{proof}
This is a finite computation over the classes $(k\bmod 2M,h\bmod M)$, carried out in exact rational arithmetic (Appendix~\ref{app:code}).
For each class we take the smallest admissible $k$ other than the dominant one.
For instance, for $S_0$ with $k'$ odd and $h'\equiv2,6\pmod 8$ the theta quotient has $k'^2E=1$; this is the dominant cusp when $k'=1$ (that is $k=4$), and otherwise $k'\ge3$ gives $E\le1/9$.
At the dominant cusps the Appell--Lerch terms used have exponent $0$.
\end{proof}

%=====================================================================
\section{Main terms}\label{sec:main}

Let $P_f$ denote the theta quotient that carries the main term:
\[
P_{S_0}=\frac{q\,\Theta(-q;q^8)\Theta(q^2;q^8)^2}{\Theta(q^3;q^8)^2},\qquad
P_{S_1}=\frac{\Theta(-q^3;q^8)\Theta(q^2;q^8)^2}{q\,\Theta(q;q^8)^2},\qquad
P_\mu=P_{U_1}=P_{U_0}=\THq.
\]
Correspondingly $f=\kappa_fP_f+(\text{Appell--Lerch terms})$, with $\kappa_f=1$ for $S_0,S_1,U_0$, $\kappa_\mu=-1$ and $\kappa_{U_1}=\frac12$; for $U_0$ the second summand is $-2U_1$.

\begin{lemma}[Poisson expansion]\label{lem:poisson}
Let $h/k$ be a cusp, $t>0$, $q=\e(h/k)e^{-t}$, and $b=a-M/2$.
Then
\[
\Theta(sq^a;q^M)=\frac1{2k}\sqrt{\frac{2\pi}{Mt}}\,e^{tb^2/(2M)}\sum_{\nu\in\Z}G_\nu\,\e\Big(\frac{\nu b}{2kM}\Big)\exp\Big(-\frac{\pi^2\nu^2}{2k^2Mt}\Big),
\]
where $G_\nu=\sum_{j=0}^{2k-1}(-s)^j\,\e\big(h(Mj(j-1)/2+aj)/k+\nu j/(2k)\big)$.
\end{lemma}

\begin{proof}
Split the sum defining $\Theta$ by $n\equiv j\pmod{2k}$.
The factor $(-s)^n\e(h(Mn(n-1)/2+an)/k)$ depends only on $j$.
Apply Poisson summation to the Gaussian $y\mapsto e^{-t(My^2/2+by)}$ along $y\in j+2k\Z$.
\end{proof}

\begin{proposition}\label{prop:main}
Let $\zeta$ be a dominant cusp of $f$ and $q=\zeta e^{-z}$, $\re z>0$.
Then
\[
P_f(\zeta e^{-z})=C_\zeta\,z^{-1/2}\,e^{\delta_f z}\,e^{a_f/z}\,F(z),
\]
with $C_\zeta,a_f,\delta_f$ as in Table~\ref{tab:const}.
Here $|F(z)-1|\le e^{\Lambda}-1$.
To define $\Lambda$, write $P_f$, via Lemmas~\ref{lem:siegel} and~\ref{lem:transf}, as a product of factors $(1-u)^{\pm1}$ coming from the products in Lemma~\ref{lem:bounds}, each with $|u|=\rho_j^{e}$ for some $e>0$.
Then $\Lambda$ is the sum of $-\log(1-|u|)$ over all these factors, counted with multiplicity.
Explicitly, $\Lambda=\sum_j|p_j|\lambda(\Ahat_j,\rho_j)+\sum_i|r_i|\lambda_\eta(\rho_i)$, where $\lambda(\Ahat,\rho)=-\log(1-\rho^{\Ahat})-\sum_{n\ge1}\big(\log(1-\rho^{n+\Ahat})+\log(1-\rho^{n-\Ahat})\big)$ (the first term omitted if $\Ahat=0$) and $\lambda_\eta(\rho)=-\sum_{n\ge1}\log(1-\rho^n)$.
The first sum runs over the theta factors and the second over the distinct eta arguments $T'$ with total multiplicity $r_i$.
The values at the conjugate cusps are the complex conjugates.
\end{proposition}

\begin{table}[ht]
\centering
\renewcommand{\arraystretch}{1.25}
\begin{tabular}{@{}lllll@{}}
\toprule
$P$ & $\zeta$ & $C_\zeta$ & $a$ & $\delta$\\
\midrule
$q\,\Theta(-q;q^8)\Theta(q^2;q^8)^2/\Theta(q^3;q^8)^2$ & $i$ & $\frac{\sqrt\pi}2e^{9\pi i/16}$ & $\pi^2/64$ & $-1/16$\\
$q^{-1}\Theta(-q^3;q^8)\Theta(q^2;q^8)^2/\Theta(q;q^8)^2$ & $i$ & $\frac{\sqrt\pi}2e^{-3\pi i/16}$ & $\pi^2/64$ & $7/16$\\
$\THq$ & $i$ & $\sqrt{\pi/2}\,e^{3\pi i/8}$ & $\pi^2/32$ & $-1/8$\\
$\THq$ & $\e(1/8)$ & $\frac{\sqrt\pi}2e^{5\pi i/16}$ & $\pi^2/128$ & $-1/8$\\
$\THq$ & $\e(3/8)$ & $\frac{\sqrt\pi}2e^{-13\pi i/16}$ & $\pi^2/128$ & $-1/8$\\
\bottomrule
\end{tabular}
\caption{Constants of the main terms.}\label{tab:const}
\end{table}

\begin{proof}
Apply Lemma~\ref{lem:siegel} and Lemma~\ref{lem:transf} (without absolute values, i.e.\ with the multipliers) to every factor.
The product of the leading powers $\e(T')^{B_2(\Ahat)/2}$ and $\e(T')^{1/24}$ gives $e^{a/z}$.
The powers of $q$ from Lemma~\ref{lem:siegel}, together with the explicit $q^{\pm1}$, give $e^{\delta z}$.
The factors $(k'M\tau-h')^{-1/2}$ give $z^{-1/2}$ and the modulus $|C_\zeta|$.
What remains is $F(z)=\prod_j\prod(1-u)^{\pm p_j}$, a product over the factors of Lemma~\ref{lem:bounds} with $|u|$ equal to powers $\rho_j^{e}$, $e>0$.
Expanding $|F-1|$ termwise gives the bound.

The phase of $C_\zeta$ is obtained from Lemma~\ref{lem:poisson}.
As $t\to0^+$ the sum is dominated by the frequencies $\pm\nu^*$ of smallest modulus with nonzero coefficient.
The argument of $C_\zeta$ is therefore $\sum_jp_j\arg\big(\sum_{\nu=\pm\nu^*_j}G_\nu\,\e(\nu b_j/(2kM_j))\big)$ plus the argument of the explicit power of $\zeta$.
Each inner sum is a finite sum of roots of unity, which we evaluate exactly; the results are the stated rational multiples of $\pi$.
For $S_1$, for instance, $\zeta^{-1}=-i$ contributes $-\pi/2$, $\Theta(-q^3;q^8)$ contributes $-\pi/16$ and $\Theta(q;q^8)^{-2}$ contributes $+3\pi/8$, for a total of $-3\pi/16$.
\end{proof}

At the dominant cusps the Appell--Lerch parts are bounded.
For $S_0$ and $S_1$, form A at $q=\pm i$ has $\Theta(-1;q^8)$ with exponent $0$, and every Lerch denominator $1-q^{8r-5}$, resp.\ $1-q^{8r-7}$, equals $1\pm i\cdot(\text{real})$ in modulus at least~$1$.
For $\mu$ and $U_1$, the form $U_0=2m(-q,q^4,-1)$ has exponent $0$ at $\pm i$.
For $U_0$, the single-$m$ forms of $U_1$ have exponent $0$ at the eighth roots of unity.

\begin{lemma}\label{lem:bessel}
For $a,N,t_0>0$,
\[
\frac1{2\pi i}\int_{t_0-i\infty}^{t_0+i\infty}z^{-1/2}e^{Nz+a/z}\,dz=\frac{\cosh(2\sqrt{aN})}{\sqrt{\pi N}}.
\]
\end{lemma}

\begin{proof}
This is the inverse Laplace transform $\mathcal L^{-1}[s^{-1/2}e^{a/s}](N)=(N/a)^{-1/4}I_{-1/2}(2\sqrt{aN})$, combined with $I_{-1/2}(x)=\sqrt{2/(\pi x)}\cosh x$.
\end{proof}

\begin{definition}[Main terms]
With $N=n+\delta_f$ define
\begin{align*}
M_{S_0}(n)&=\cos\Big(\frac{9\pi}{16}-\frac{\pi n}2\Big)\frac{\cosh(\frac\pi4\sqrt N)}{\sqrt N}, \\
M_{S_1}(n)&=\cos\Big(\frac{3\pi}{16}+\frac{\pi n}2\Big)\frac{\cosh(\frac\pi4\sqrt N)}{\sqrt N},\\
M_{\mu}(n)&=\sqrt2\cos\Big(\frac{5\pi}8+\frac{\pi n}2\Big)\frac{\cosh(\frac{\pi}{2\sqrt2}\sqrt N)}{\sqrt N}, \\
M_{U_1}(n)&=\frac1{\sqrt2}\cos\Big(\frac{3\pi}8-\frac{\pi n}2\Big)\frac{\cosh(\frac{\pi}{2\sqrt2}\sqrt N)}{\sqrt N},\\
M_{U_0}(n)&=\Big[\cos\Big(\frac{5\pi}{16}-\frac{\pi n}4\Big)+\cos\Big(\frac{13\pi}{16}+\frac{3\pi n}4\Big)\Big]\frac{\cosh(\frac{\pi}{4\sqrt2}\sqrt N)}{\sqrt N}.
\end{align*}
\end{definition}

Each $M_f(n)$ is $\sum_\zeta\zeta^{-n}C_\zeta\cosh(2\sqrt{a_fN})/\sqrt{\pi N}$, summed over the dominant cusps (Lemma~\ref{lem:bessel}), multiplied by $\kappa_f$.

\begin{lemma}\label{lem:trig}
Let $\theta_f(n)$ denote the trigonometric factor of $M_f(n)$, including its constant.
Then $|\theta_f(n)|\ge L_f$ for all $n$, where $L_{S_0}=\cos\frac{7\pi}{16}$, $L_{S_1}=\cos\frac{5\pi}{16}$, $L_\mu=\sqrt2\cos\frac{3\pi}8$, $L_{U_1}=\frac1{\sqrt2}\cos\frac{3\pi}8$, and $L_{U_0}=\min_{0\le r<8}|\theta_{U_0}(r)|>0.2758$.
Moreover $\operatorname{sign}\theta_f(n)$ is the sign listed in Table~\ref{tab:main} for $n$.
\end{lemma}

\begin{proof}
$\theta_f(n)$ depends only on $n\bmod 4$ (resp.\ $8$); evaluate the finitely many values.
\end{proof}

\begin{theorem}\label{thm:asym}
For each $f$, $c_f(n)=M_f(n)+O\big(n^{3/4}\exp(\tfrac\pi4\sqrt{\varepsilon_f n})\big)$, with $\varepsilon_f$ from Table~\ref{tab:exp}.
\end{theorem}

This follows from Theorem~\ref{thm:explicit} and Lemma~\ref{lem:large}.

%=====================================================================
\section{The explicit circle method}\label{sec:circle}

Fix $f$ and $n\ge1$.
Put $t_0=\frac\pi8\sqrt{\varepsilon_f/n}$ and $K=\lfloor\sqrt{2\pi/t_0}\rfloor$.
By Cauchy's formula with $q=e^{-w}$,
\[
c_f(n)=\frac1{2\pi i}\int_{t_0-\pi i}^{t_0+\pi i}f(e^{-w})\,e^{nw}\,dw.
\]
Dissect the segment by the Farey fractions $h/k$ of order $K$.
The arc of $h/k$ consists of the $w=z-2\pi ih/k$ with $z=t_0-i\varphi$ and $-\vartheta_1\le\varphi\le\vartheta_2$, where $\vartheta_{1,2}=2\pi/(k(k+k_{1,2}))$ and $k_1,k_2$ are the denominators of the neighbouring fractions.

\begin{lemma}[Geometry]\label{lem:geom}
On every arc, $t_0\le k|z|<\sqrt{4\pi t_0}$ and $\re(1/z)/k^2>1/(4\pi)$.
Consequently every $\rho$ in Lemma~\ref{lem:transf} satisfies $\rho\le\exp(-\pi g^2/M)$.
On the arc of a cusp of order $k$ we have $\re(1/z)\ge1/(t_0+2\pi/k^2)$, and the arc extends at least $Y_k=2\pi/\big(k(\sqrt{2\pi/t_0}+k)\big)$ on both sides.
\end{lemma}

\begin{proof}
Since $k_i\le K$ and $k+k_i\ge K+1$, we have $2\pi/(k(k+K))\le\vartheta_i\le2\pi/(k(K+1))$.
Then $k^2|z|^2=k^2t_0^2+k^2\varphi^2\le K^2t_0^2+(2\pi/(K+1))^2<2\pi t_0+2\pi t_0$, and $\re(1/z)=t_0/|z|^2$.
\end{proof}

\begin{lemma}[Lerch numerators]\label{lem:lerch}
Let $Q=q^M$ with $|q|=e^{-t_0}$, and let
\[
N(x,z)=\sum_{r\in\Z}\frac{(-1)^rQ^{r(r-1)/2}z^r}{1-Q^{r-1}xz}
\]
be the sum in \eqref{eq:m}.
For the representations of Proposition~\ref{prop:exponents},
\[
|N(x,z)|\le e^{c\,t_0}\,\frac{1+\sqrt{\pi/(\kappa t_0)}}{1-e^{-d\,t_0}},
\]
with $\kappa=4$ for $M=8$, $\kappa=2$ for $M=4$, and $(c,d)$ given below; here $d$ is the least $|e|$ among the denominators $1\pm q^{e}$.
\[
\renewcommand{\arraystretch}{1.15}
\begin{array}{@{}lllcc@{}}
\toprule
f & \text{form} & \text{denominators} & d & c\\ \midrule
S_0 & m(-q^3,q^8,-1) & 1-q^{8r-5} & 3 & 1\\
S_0 & m(-q^3,q^8,-q^2) & 1-q^{8r-3} & 3 & \tfrac14\\
S_0 & m(-q^3,q^8,-q^6) & 1-q^{8r+1} & 1 & \tfrac14\\
S_1 & m(-q,q^8,-1) & 1-q^{8r-7} & 1 & 1\\
S_1 & m(-q,q^8,-q^2) & 1-q^{8r-5} & 3 & \tfrac14\\
S_1 & m(-q,q^8,-q^6) & 1-q^{8r-1} & 1 & \tfrac14\\
U_0 & m(-q,q^4,-1) & 1-q^{4r-3} & 1 & \tfrac12\\
U_0 & m(-q,q^4,q^3) & 1+q^{4r} & 4 & \tfrac18\\
U_1 & m(-q,q^4,-q^2) & 1-q^{4r-1} & 1 & 0\\
U_1 & m(-q,q^4,q) & 1+q^{4r-2} & 2 & \tfrac18\\ \bottomrule
\end{array}
\]
\end{lemma}

\begin{proof}
For either sign of $e\neq0$, $|1\pm q^e|\ge1-e^{-|e|t_0}$; for $U_0$ with $z=q^3$ the single term $r=0$ has denominator exactly $2$.
The numerators satisfy $|Q^{r(r-1)/2}z^r|=e^{-\kappa t_0(r-r_0)^2+ct_0}$ after completing the square, and $\sum_{r\in\Z}e^{-\kappa t_0(r-r_0)^2}\le1+\sqrt{\pi/(\kappa t_0)}$.
\end{proof}

\begin{lemma}[Tails]\label{lem:tail}
For $a,N>0$, $t_0>0$ and $Y>0$,
\[
\Big|\frac1{2\pi i}\int_{t_0+iY}^{t_0+i\infty}z^{-1/2}e^{Nz+a/z}\,dz\Big|\le\frac{e^{Nt_0}}{2\pi N}\,e^{at_0/Y^2}\Big(2Y^{-1/2}+\frac{2a}3Y^{-3/2}\Big),
\]
and the same bound holds for the integral from $t_0-i\infty$ to $t_0-iY$.
\end{lemma}

\begin{proof}
Integrate by parts once, using $|e^{Nz}|=e^{Nt_0}$, $|e^{a/z}|\le e^{at_0/Y^2}$ for $|\im z|\ge Y$, and $|\frac{d}{dz}(z^{-1/2}e^{a/z})|\le(\frac12|z|^{-3/2}+a|z|^{-5/2})e^{at_0/Y^2}$.
\end{proof}

\begin{theorem}[Explicit error bound]\label{thm:explicit}
For every $f$ and $n\ge1$, $|c_f(n)-M_f(n)|\le B_f(n)$, where
\[
B_f(n)=e^{nt_0}\Big[C_f(t_0)\,e^{\pi^2\varepsilon_f/(64t_0)}+D_f(t_0)\Big]+T_f(n),
\]
and $C_f,D_f,T_f$ are as follows.
\begin{enumerate}
\item $C_f(t_0)$ is the maximum over the cusp classes of the sum, over the terms of the chosen representation, of the product of:
the bound of Lemma~\ref{lem:lerch} (for Appell--Lerch terms);
the factor $\prod_j(2\pi g_j/M_j)^{p_j/2}(k|z|)^{-\sum p_j/2}$ from Lemma~\ref{lem:transf}(1), with $k|z|$ replaced by $t_0$ or $\sqrt{4\pi t_0}$ according to the sign of the exponent;
the $|q|$-powers of Lemma~\ref{lem:siegel} and the explicit powers $q^{\pm1}$ of the representation (only factors larger than $1$ are kept);
and the products $U^{p_j}$ ($p_j>0$) or $L^{p_j}$ ($p_j<0$) of Lemma~\ref{lem:bounds} at $\rho_j=\exp(-\pi g_j^2/M_j)$.
\item $D_f(t_0)$ bounds the integrand minus the main term on the dominant arcs. With $X=1/(t_0+2\pi/k^2)$ it is
\[
|\kappa_f|\,|C_\zeta|\,t_0^{-1/2}e^{\max(\delta_f,0)t_0}\,e^{a_fX}\big(e^{\Lambda}-1\big),
\]
with $\Lambda$ evaluated at $\re(1/z)=X$, plus the bound of item~(1) for the Appell--Lerch terms, evaluated with $\re(1/z)\ge X$.
\item $T_f(n)$ is twice the number of dominant cusps, times $|\kappa_f|\,|C_\zeta|$, times the bound of Lemma~\ref{lem:tail} with $a=a_f$, $N=n+\delta_f$ and $Y=Y_k$.
\end{enumerate}
\end{theorem}

\begin{proof}
The Farey arcs partition the segment of integration, which has length $2\pi$, so a pointwise bound for the integrand on each arc bounds the integral.
On a non-dominant arc, the chosen representation and Lemmas~\ref{lem:siegel}--\ref{lem:bounds}, \ref{lem:geom} and \ref{lem:lerch} bound $|f|$ by $C_f(t_0)\exp(\frac{\pi^2}{64}E\,\re(1/z))$.
Here $E\le\varepsilon_f$ by Proposition~\ref{prop:exponents}, and $\re(1/z)\le1/t_0$; also $|e^{nw}|=e^{nt_0}$.
On a dominant arc we write $f=\pm(\text{main term})+(P_f-\text{main term})+(\text{Appell--Lerch terms})$ and use Proposition~\ref{prop:main}.
The main term integrated over the arc equals the full vertical integral of Lemma~\ref{lem:bessel} minus two tails, which Lemma~\ref{lem:tail} bounds with $Y=Y_k$ (Lemma~\ref{lem:geom}).
Summing the full integrals over the dominant cusps gives $M_f(n)$.
\end{proof}

Since $nt_0+\pi^2\varepsilon_f/(64t_0)=\frac\pi4\sqrt{\varepsilon_fn}$, the bound grows like $\exp(\frac\pi4\sqrt{\varepsilon_fn})$ times a power of $n$.
The main term grows like $\exp(2\sqrt{a_fn})=\exp(\frac\pi4\sqrt{\alpha_fn})$, where $\alpha_f=64a_f/\pi^2$ is the dominant exponent of Table~\ref{tab:exp} and $\alpha_f>\varepsilon_f$.

%=====================================================================
\section{Proof of the Main Theorem}\label{sec:proof}

\subsection{Certified evaluation}
The constants in $B_f(n)$ are explicit convergent products of the form $U$, $L$, $\prod(1\pm\rho^n)$, and the sums $\Lambda$, at finitely many values of $\rho$.
We evaluate the products in $50$-digit decimal arithmetic with explicit truncation bounds ($\log(1+x)\le x$, and $-\log(1-x)\le2x$ for $0\le x\le\frac12$).
We round outward, and replace each $\rho$ by an upper bound on a grid of mesh $10^{-7}$; this is conservative because $U$ increases and $L$ decreases in $\rho$.
The remaining elementary operations are performed in IEEE double precision, with relative error below $10^{-14}$.
We therefore require a logarithmic margin of at least $10^{-8}$.

\begin{proposition}\label{prop:cert}
For each $f$ and every integer $n$ with $N_0(f)\le n\le 2\cdot10^5$,
\[
\log\Big(L_f\,\frac{\cosh(2\sqrt{a_fN})}{\sqrt N}\Big)-\log B_f(n)>10^{-8},\qquad N=n+\delta_f,
\]
where $N_0(f)$ is given in Table~\ref{tab:N0}.
Moreover this margin is increasing in $n$ on $[N_0(f),2\cdot10^5]$.
\end{proposition}

\begin{table}[ht]
\centering
\renewcommand{\arraystretch}{1.2}
\begin{tabular}{@{}lrcc@{}}
\toprule
$f$ & $N_0(f)$ & margin at $N_0$ & margin at $2\cdot10^5$\\
\midrule
$S_0$ & $11231$ & $1.0\cdot10^{-4}$ & $131.8$\\
$S_1$ & $10652$ & $4.3\cdot10^{-4}$ & $132.9$\\
$\mu$ & $976$ & $7.3\cdot10^{-4}$ & $225.1$\\
$U_1$ & $493$ & $1.0\cdot10^{-2}$ & $308.1$\\
$U_0$ & $5180$ & $1.7\cdot10^{-3}$ & $100.1$\\
\bottomrule
\end{tabular}
\caption{Certified thresholds and logarithmic margins.}\label{tab:N0}
\end{table}

\begin{lemma}\label{lem:large}
The inequality of Proposition~\ref{prop:cert} holds for all $n\ge2\cdot10^5$.
\end{lemma}

\begin{proof}
Write $t=t_0(n)$, a decreasing function of $n$ proportional to $n^{-1/2}$, and $n^*=2\cdot10^5$, $t^*=t_0(n^*)$.
Every factor $\phi(t)$ occurring in $C_f$, $D_f$ and $T_f$ satisfies: $t^{\gamma}\phi(t)$ is nondecreasing on $(0,t^*]$ for some $\gamma\le\frac32$.
The factors are of three kinds:
\begin{itemize}
\item $t^{-1/2}$, $(\sqrt{4\pi t})^{1/2}$, $e^{ct}$, $1+\sqrt{\pi/(\kappa t)}$ (note $t^{1/2}(1+\sqrt{\pi/(\kappa t)})=t^{1/2}+\sqrt{\pi/\kappa}$), and $(1-e^{-dt})^{-1}$ (note $t/(1-e^{-dt})$ is increasing);
\item the major-arc factor $e^{a_fX}(e^{\Lambda}-1)$, which is nondecreasing in $\rho$ and hence in $t$, being a series in powers of $\rho$ whose exponents all exceed that of $e^{a_fX}$;
\item the tail factor, which is polynomial in $Y_k\asymp t^{1/2}$.
\end{itemize}
Hence $B_f(n)\le B_f(n^*)\,(n/n^*)^{3/4}\exp\big(\frac\pi4\sqrt{\varepsilon_f}(\sqrt n-\sqrt{n^*})\big)$ for $n\ge n^*$.
The logarithm of the main-term lower bound has derivative at least $\frac{\pi\sqrt{\alpha_f}}{8\sqrt N}-\frac1{2N}$, while the logarithm of this upper bound has derivative at most $\frac{\pi\sqrt{\varepsilon_f}}{8\sqrt n}+\frac3{4n}$.
Since $\alpha_f>\varepsilon_f$, the margin is increasing once $\sqrt n>10/(\pi(\sqrt{\alpha_f}-\sqrt{\varepsilon_f}))$, which holds for $n\ge n^*$.
At $n^*$ the margin is positive by Proposition~\ref{prop:cert}.
\end{proof}

\begin{proof}[Proof of the Main Theorem]
Let $n\ge N_0(f)$.
By Theorem~\ref{thm:explicit}, Proposition~\ref{prop:cert} and Lemma~\ref{lem:large},
$|c_f(n)-M_f(n)|\le B_f(n)<L_f\cosh(2\sqrt{a_fN})/\sqrt N\le|M_f(n)|$.
So $c_f(n)$ is nonzero and has the sign of $M_f(n)$, which by Lemma~\ref{lem:trig} is the sign in Table~\ref{tab:main}.
For $0\le n<N_0(f)$ we computed $c_f(n)$ exactly from the defining series \eqref{eq:S}--\eqref{eq:mu} in integer arithmetic; the signs and the exceptions are as in Table~\ref{tab:main}.
\end{proof}

\begin{remark}
The bounds on the dominant arcs are computed at one cusp of each conjugate pair.
This suffices because $q\mapsto\bar q$ maps the arc of $h/k$ onto that of $(k-h)/k$ and preserves all moduli.
\end{remark}

%=====================================================================
\section{Further remarks}\label{sec:remarks}

\subsection*{Third-order functions}
Dastidar~\cite{Da} conjectured sign laws for the third-order functions $\phi(q)=\sum q^{n^2}/(-q^2;q^2)_n$ and $\chi(q)=\sum q^{n^2}(-q;q)_n/(-q^3;q^3)_n$.
Computing their coefficients for $n\le1500$, we find that $\phi$ has signs $+,+,-,-$ for $n\equiv0,1,2,3\pmod4$, with zeros only at $n=2,8,10,26$.
We also find that $\chi$ has signs $+,+,+,-,-,-$ for $n\equiv0,\dots,5\pmod 6$, with zeros exactly at $n=3,4,5,8,9,11,16,17,18,20,21,23,27,32,35,41,44,47,68,71,80$.
The method of this paper should apply to both.

\subsection*{Period and dominant cusps}
In all cases considered here, and in Dastidar's theorem for $\rho(q)$, the period of the sign law equals the order of the dominant roots of unity.
These roots form a single Galois orbit at which a pure theta quotient dominates and the Appell--Lerch part is bounded.
It would be interesting to determine which classical mock theta functions have this property, and so to classify their periodic sign laws.

%=====================================================================
\appendix
\sloppy
\section{Computations}\label{app:code}

The programs needed for the proof (Python~3, standard library only) accompany this paper as ancillary files and are available at \url{https://github.com/ramanujamath/mock-theta-sign-laws}, in the directory \href{https://github.com/ramanujamath/mock-theta-sign-laws/tree/main/code}{\texttt{code}}.
\begin{itemize}
\item \href{https://github.com/ramanujamath/mock-theta-sign-laws/blob/main/code/coeffs.py}{\texttt{coeffs.py}}: exact integer coefficients of $S_0,S_1,\mu,U_0,U_1$ from \eqref{eq:S}--\eqref{eq:mu} for $n\le12000$, and the check of Table~\ref{tab:main} in that range (used in the proof of the Main Theorem for $n<N_0$).
\item \href{https://github.com/ramanujamath/mock-theta-sign-laws/blob/main/code/exponents_S.py}{\texttt{exponents\_S.py}}: the exponent classes for $S_0,S_1$ (Proposition~\ref{prop:exponents}).
\item \href{https://github.com/ramanujamath/mock-theta-sign-laws/blob/main/code/exact_phase.py}{\texttt{exact\_phase.py}}: the exact phases of $C_\zeta$ (Proposition~\ref{prop:main}).
\item \href{https://github.com/ramanujamath/mock-theta-sign-laws/blob/main/code/farey_lib.py}{\texttt{farey\_lib.py}}, \href{https://github.com/ramanujamath/mock-theta-sign-laws/blob/main/code/farey_bound.py}{\texttt{farey\_bound.py}}, \href{https://github.com/ramanujamath/mock-theta-sign-laws/blob/main/code/farey_N0.py}{\texttt{farey\_N0.py}}, \href{https://github.com/ramanujamath/mock-theta-sign-laws/blob/main/code/farey_S1_N0.py}{\texttt{farey\_S1\_N0.py}}, \href{https://github.com/ramanujamath/mock-theta-sign-laws/blob/main/code/farey_gen.py}{\texttt{farey\_gen.py}}: the bound $B_f(n)$ of Theorem~\ref{thm:explicit}. The last also computes the exponent classes for $\mu,U_1,U_0$.
\item \href{https://github.com/ramanujamath/mock-theta-sign-laws/blob/main/code/certify.py}{\texttt{certify.py}}: outward-rounded multiprecision evaluation of the products of Lemma~\ref{lem:bounds}.
\item \href{https://github.com/ramanujamath/mock-theta-sign-laws/blob/main/code/certify_run.py}{\texttt{certify\_run.py}}: the verification of Proposition~\ref{prop:cert} and the sparse values used in Lemma~\ref{lem:large}.
\end{itemize}
The total running time is a few minutes on a laptop.

%=====================================================================
\begin{thebibliography}{99}

\bibitem{An} G.~E. Andrews, \emph{On the theorems of Watson and Dragonette for Ramanujan's mock theta functions}, Amer.\ J. Math.\ \textbf{88} (1966), 454--490.

\bibitem{Ap} T.~M. Apostol, \emph{Modular Functions and Dirichlet Series in Number Theory}, 2nd ed., Graduate Texts in Mathematics~41, Springer, New York, 1990.

\bibitem{BO} K.~Bringmann and K.~Ono, \emph{The $f(q)$ mock theta function conjecture and partition ranks}, Invent.\ Math.\ \textbf{165} (2006), 243--266.

\bibitem{Da} M.~G. Dastidar, \emph{Sign laws and mock theta functions}, arXiv:2606.31606 (2026).

\bibitem{Dr} L.~A. Dragonette, \emph{Some asymptotic formulae for the mock theta series of Ramanujan}, Trans.\ Amer.\ Math.\ Soc.\ \textbf{72} (1952), 474--500.

\bibitem{GM} B.~Gordon and R.~J. McIntosh, \emph{Some eighth order mock theta functions}, J. London Math.\ Soc.\ (2) \textbf{62} (2000), 321--335.

\bibitem{HM} D.~Hickerson and E.~Mortenson, \emph{Hecke-type double sums, Appell--Lerch sums, and mock theta functions, I}, Proc.\ London Math.\ Soc.\ (3) \textbf{109} (2014), 382--422; arXiv:1208.1421.

\bibitem{KL} D.~S. Kubert and S.~Lang, \emph{Modular Units}, Grundlehren der Mathematischen Wissenschaften~244, Springer, New York, 1981.

\bibitem{Mc} R.~J. McIntosh, \emph{Second order mock theta functions}, Canad.\ Math.\ Bull.\ \textbf{50} (2007), 284--290.

\end{thebibliography}

\end{document}
